\documentclass[12pt,a4paper]{article}

\usepackage{amsmath, amsthm, amssymb, amsfonts}
\usepackage{mathrsfs}
\usepackage{hyperref}
\usepackage{geometry}
\usepackage{enumitem}
\usepackage{booktabs}

\geometry{margin=2.5cm}

\newtheorem{theorem}{Theorem}[section]
\newtheorem{proposition}[theorem]{Proposition}
\newtheorem{lemma}[theorem]{Lemma}
\newtheorem{corollary}[theorem]{Corollary}
\newtheorem{definition}[theorem]{Definition}
\newtheorem{remark}[theorem]{Remark}
\newtheorem{assumption}[theorem]{Assumption}
\newtheorem{openproblem}[theorem]{Open Problem}
\newtheorem{conjecture}[theorem]{Conjecture}

\newcommand{\Ai}{\operatorname{Ai}}
\newcommand{\lam}[1]{\lambda_{#1}^{(c)}}
\newcommand{\eps}[1]{\varepsilon_{#1}^{(c)}}
\newcommand{\pswf}[1]{\psi_{#1}^{(c)}}
\newcommand{\norm}[1]{\left\|#1\right\|}
\newcommand{\Tc}{T_c}
\newcommand{\VN}{V_N}
\newcommand{\Dc}{D_c}

\title{%
  \textbf{Paper~XI: Fredholm Determinant Microlocal Inversion}\\[0.4em]
  \large Spectral Gap at the Galerkin Scale via
  $\log \det(I - zT_c)$
}
\author{\textit{prolate-gram-coercivity programme}}
\date{May 2026}

\begin{document}
\maketitle

\begin{abstract}
Paper~X established that the best available lower bound
on $\varepsilon_N^{(c)} = 1 - \lam{N}$ from local
(WKB/Pr\"ufer/ORX) methods is
$\varepsilon_N \gtrsim e^{-Cc^{2/3}}$,
which falls short of the required polynomial bound
$\varepsilon_N \geq Cc^{-1/2}$ by a superpolynomial factor.
The gap reflects a fundamental structural mismatch:
local methods describe individual eigenfunctions,
but the target bound at the Galerkin scale
$k = O(c^{1/3}) \ll N_{\mathrm{Sh}}$
is an \emph{intermediate spectral problem}
inaccessible to purely local analysis.

\medskip\noindent
This paper develops the Fredholm determinant approach
as the unique structural route capable of
bridging global spectral information
to a single eigenvalue at small index.
We introduce the Fredholm determinant
$\Dc(z) = \det(I - z\Tc)$,
derive its asymptotic expansion in the bulk-scaled
regime $z \approx 1$,
and identify the \emph{missing bridge} between
$\log \Dc(z)$ and $\lam{N}$.

\medskip\noindent
We also confront the fundamental tension:
$\varepsilon_N \geq c^{-1/2}$ requires the prolate
operator to \emph{leak energy} from superconcentrated
modes at a polynomial rate, which is \emph{counter to}
the natural concentration structure of PSWFs.
We make this tension precise,
and formulate the conditions under which it resolves
in favour of B-strong.
\end{abstract}

\tableofcontents
\bigskip\hrule\bigskip

%% ============================================================
\section{The Structural Tension}
\label{sec:tension}
%% ============================================================

\subsection{What B-strong requires against what PSWFs do}

The PSWF concentration eigenvalues $\lam{k}$
are classically described as follows:
\begin{itemize}[nosep]
  \item For $k \leq N_{\mathrm{Sh}} = 2c/\pi$:
        $\lam{k}$ is exponentially close to $1$;
        the modes $\pswf{k}$ are \emph{maximally concentrated}
        in $[-1,1]$.
  \item The transition from $\lam{k} \approx 1$ to
        $\lam{k} \approx 0$ occurs on a window
        $|k - N_{\mathrm{Sh}}| = O(c^{1/3})$.
  \item For $k \ll N_{\mathrm{Sh}}$: the concentration
        is \emph{superexponentially stable};
        the energy leak $\varepsilon_k = 1 - \lam{k}$
        is of order $e^{-c^\alpha}$ for some $\alpha > 0$.
\end{itemize}

\begin{remark}[The tension, made explicit]
\label{rem:tension}
B-strong requires $\varepsilon_N \geq Cc^{-1/2}$,
i.e.\ the prolate operator must lose at least
a polynomial fraction $c^{-1/2}$ of energy
for the $N$-th mode.
But the entire structure of PSWF theory says that
for $k \leq N = O(c^{1/3}) \ll N_{\mathrm{Sh}}$,
the energy loss is superexponentially \emph{smaller}
than any polynomial.

\medskip\noindent
This is not a technical failure --- it is a
\emph{structural contradiction} between
the known superconcentration asymptotics and
the B-strong requirement.
One of two things must therefore be true:
\begin{enumerate}[nosep]
  \item[(A)] The known asymptotics are insufficiently
        sharp at the scale $k = O(c^{1/3})$,
        and the true $\varepsilon_N$ is in fact
        polynomial (not exponential).
  \item[(B)] B-strong is false as stated,
        i.e.\ $\varepsilon_N \ll c^{-1/2}$,
        and the programme must either weaken
        the assumption or seek an alternative
        route to coercivity.
\end{enumerate}
\noindent
\textbf{Resolving this dichotomy is the central
problem of Paper~XI.}
\end{remark}

\subsection{Why local methods cannot resolve the tension}

Local methods (WKB, ORX, Landau--Pollak) give information of
the form:
\[
  \varepsilon_k \sim \exp(-c \cdot \phi(k/c))
\]
with $\phi(t) = \pi(1 - 2t/\pi) > 0$ for $t < \pi/2$.
For $k = O(c^{1/3})$: $k/c = O(c^{-2/3}) \to 0$,
so $\phi(k/c) \to \pi > 0$, giving
$\varepsilon_k \sim e^{-\pi c}$ --- superexponentially small.

This estimate cannot be improved by any refinement
of the same local asymptotic:
the exponential suppression comes from tunneling
through the effective potential barrier
$V_{\mathrm{eff}}(x) = c^2(x^2-1)$ at $|x| > 1$,
which is a genuine physical mechanism, not an artifact.

\begin{proposition}[Local methods are optimally wrong]
\label{prop:localfail}
Any estimate of the form
$\varepsilon_k \geq f(c)$ derived solely from the
WKB outer-tail of $\pswf{k}$ (Proposition~3.1 of Paper~X)
will satisfy $f(c) \leq e^{-c^{\delta}}$ for some
$\delta > 0$, and hence $f(c) \ll c^{-1/2}$.
Local analysis cannot distinguish scenario~(A)
from scenario~(B) of Remark~\ref{rem:tension}.
\end{proposition}

%% ============================================================
\section{The Fredholm Determinant}
\label{sec:fredholm}
%% ============================================================

\subsection{Definition and basic properties}

\begin{definition}[Prolate Fredholm determinant]
For $z \in \mathbb{C}$ with $|z| \leq 1$:
\[
  \Dc(z) := \det(I - z\Tc)
  = \prod_{k=0}^{\infty} (1 - z\lam{k}).
\]
Since $\Tc$ is trace-class (as an integral operator
on $L^2([-1,1])$ with continuous kernel),
this product converges absolutely.
\end{definition}

\begin{proposition}[Basic identities]
\label{prop:fredholm_basic}
\begin{enumerate}[nosep]
  \item $\Dc(1) = \prod_k (1 - \lam{k}) = \prod_k \eps{k}$.
  \item $\partial_z \log \Dc(z) = -\sum_k \frac{\lam{k}}{1-z\lam{k}}$.
  \item $\log \Dc(z) = -\sum_{n=1}^{\infty}
        \frac{z^n}{n}\operatorname{tr}(\Tc^n)
        = -\sum_{n=1}^{\infty} \frac{z^n}{n}
        \sum_k (\lam{k})^n$.
  \item $\operatorname{tr}(\Tc^n)
        = \int_{[-1,1]^n}
        \prod_{j=1}^n K_c(x_j, x_{j+1})\,
        dx_1\cdots dx_n$
        (with $x_{n+1} = x_1$),
        where $K_c(x,y) = \sin(c(x-y))/(\pi(x-y))$.
\end{enumerate}
\end{proposition}

\subsection{Why Fredholm reaches the Galerkin scale}

The key insight is that $\Dc(z)$ is \emph{global}:
it contains information about all eigenvalues
$\{\lam{k}\}_{k \geq 0}$ simultaneously.
In particular, $\Dc(z)$ near $z = \lam{N}^{-1}$
has a simple zero whose location encodes $\lam{N}$.

\begin{proposition}[Zero extraction]
\label{prop:zero}
For each $k$, $\Dc(z)$ has a simple zero at
$z = z_k := (\lam{k})^{-1} > 1$.
In particular, for $z$ near $z_N^{-1} = \lam{N}$
(approaching from below along the real axis):
\[
  \varepsilon_N = 1 - \lam{N}
  = 1 - z_N^{-1}
  = \frac{z_N - 1}{z_N}.
\]
\noindent
Equivalently: extracting $\varepsilon_N$ from $\Dc$
requires locating the $N$-th \emph{smallest zero}
of $z \mapsto \Dc(1/z)$ in $(0,1)$.
\end{proposition}

\subsection{The missing bridge}

\begin{definition}[Partial determinant and truncated log-derivative]
For $M \geq 0$, define the \emph{Galerkin partial determinant}:
\[
  \Dc^{(N)}(z) := \prod_{k=0}^{N} (1 - z\lam{k}),
\]
so that $\Dc(z) = \Dc^{(N)}(z) \cdot \Dc^{(>N)}(z)$
with $\Dc^{(>N)}(z) = \prod_{k > N}(1 - z\lam{k})$.
\end{definition}

\begin{remark}[The missing bridge, precisely]
\label{rem:bridge}
To extract $\varepsilon_N$ from $\Dc$, one needs
to separate $\Dc^{(N)}$ from $\Dc^{(>N)}$.
This requires:
\begin{enumerate}[nosep]
  \item \textbf{Asymptotic control of $\log \Dc(z)$}
        for $z \in (0, 1)$ in the
        \emph{bulk-scaled regime} $z = 1 - sc^{-1/2}$
        (not in the edge-scaled regime $z \approx \lam{N_{\mathrm{Sh}}}^{-1}$).
  \item \textbf{Separation of contributions:}
        the sum $\sum_k \lam{k}^n$ splits as
        $\sum_{k \leq N} \lam{k}^n + \sum_{k > N}\lam{k}^n$;
        for $k \leq N$ all terms are $\approx 1$,
        while for $k > N_{\mathrm{Sh}} + C c^{1/3}$
        all terms are $\approx 0$.
        The intermediate range
        $k \in [N, N_{\mathrm{Sh}} + Cc^{1/3}]$
        contributes the non-trivial part.
  \item \textbf{Finite-difference log-derivative:}
        the difference
        $\log\Dc^{(N)}(z) - \log\Dc^{(N-1)}(z)
        = \log(1 - z\lam{N})$
        gives $\lam{N}$ directly.
\end{enumerate}
The obstacle is step (1): the asymptotics of
$\log \Dc(z)$ in the bulk regime $z \approx 1$
are governed by $\operatorname{tr}(\Tc^n)$,
which requires controlling oscillatory integrals
of the form
$\int_{[-1,1]^n} \prod_j K_c(x_j, x_{j+1})\,dx$
for large $n$.
This is a \emph{multi-dimensional stationary phase problem}.
\end{remark}

%% ============================================================
\section{Bulk-Scaled Asymptotics of $\log D_c(z)$}
\label{sec:bulk}
%% ============================================================

\subsection{Trace asymptotics}

The standard asymptotic for the trace is:
\[
  \operatorname{tr}(\Tc) = \sum_k \lam{k}
  \sim \frac{2c}{\pi}
  = N_{\mathrm{Sh}},
\]
and more generally
$\sum_k (\lam{k})^n \sim N_{\mathrm{Sh}}$
for each fixed $n$ (since $\lam{k} \approx 1$
for $k \leq N_{\mathrm{Sh}}$ and $\approx 0$ otherwise).

\begin{proposition}[Log-determinant in bulk regime]
\label{prop:logdet}
For $z = 1 - sc^{-1/2}$ with $s > 0$ fixed:
\[
  \log \Dc(z)
  = \sum_k \log(1 - z\lam{k})
  \approx \sum_{k \leq N_{\mathrm{Sh}}} \log(1 - z)
  + \sum_{\substack{k : |\lam{k} - z^{-1}| = O(c^{-1/3})}}
  \log(1 - z\lam{k}).
\]
The first sum contributes
$N_{\mathrm{Sh}} \log(sc^{-1/2}) \sim -\frac{c}{2}\log c$,
which dominates.
The second sum (edge correction) is $O(c^{1/3}\log c)$.
The Galerkin terms $k \leq N$ contribute
$\sum_{k=0}^N \log(1 - z\lam{k})
\approx N \log(sc^{-1/2})
= O(c^{1/3}\log c)$,
subdominant relative to the full sum.
\end{proposition}

\begin{remark}[Consequence for microlocal inversion]
Proposition~\ref{prop:logdet} shows that the Galerkin
terms $k \leq N$ contribute at the level
$O(c^{1/3}\log c)$ to $\log \Dc(z)$,
which is exponentially smaller than the total
$O(c \log c)$.
Extracting $\lam{N}$ from $\log \Dc(z)$ therefore
requires a relative precision of $c^{-2/3}$ in
the log-determinant --- comparable to the Galerkin
contribution itself.
This is precisely the difficulty:
$\Dc(z)$ is globally dominated by the bulk,
but B-strong concerns the exponentially suppressed
Galerkin sector.
\end{remark}

\subsection{The required precision and its cost}

\begin{proposition}[Precision requirement]
\label{prop:precision}
To determine $\varepsilon_N = 1 - \lam{N}$
with absolute precision $\delta$ via
$\log \Dc(z)$, one requires control of
$\log \Dc(z)$ to precision $\delta/\varepsilon_N$.
Since $\varepsilon_N \sim e^{-Cc^{2/3}}$
(Paper~X, Proposition~3.1),
this requires precision $\delta \cdot e^{Cc^{2/3}}$
in $\log \Dc$ --- superexponentially better
than the current $O(c^{1/3} \log c)$ error.
\end{proposition}

\begin{corollary}[Fredholm inversion is obstructed]
\label{cor:fredholm_obstructed}
If the WKB bound $\varepsilon_N \sim e^{-Cc^{2/3}}$
is sharp (scenario~(B) of Remark~\ref{rem:tension}),
then $\log \Dc(z)$ cannot distinguish
$\varepsilon_N$ from $0$ at any accessible
asymptotic precision.
If the bound is \emph{not} sharp (scenario~(A)),
the polynomial $\varepsilon_N \sim c^{-1/2}$ must
emerge from a \emph{cancellation mechanism}
in the Fredholm determinant that reduces the
apparent exponential suppression.
\end{corollary}

%% ============================================================
\section{The Dichotomy: Scenario (A) vs.\ (B)}
\label{sec:dichotomy}
%% ============================================================

\subsection{Scenario (A): polynomial gap despite superconcentration}

For scenario~(A) to hold, there must exist a mechanism
that generates \emph{polynomial energy leakage}
$\varepsilon_N \sim c^{-1/2}$ despite the exponential
suppression visible in WKB.
Three candidate mechanisms:

\begin{enumerate}
  \item \textbf{Collective eigenvalue interaction.}
        The prolate eigenvalues are not independent;
        they repel each other (as in GUE).
        Eigenvalue repulsion at the Galerkin scale
        could force $\varepsilon_N$ to be larger
        than the individual WKB estimate suggests.
        Specifically, the product
        $\prod_{k=0}^N \varepsilon_k$ is controlled
        by $\Dc(1)$, and if $\Dc(1)$ has polynomial
        order, then at least one $\varepsilon_k$ must
        be polynomial.

  \item \textbf{WKB estimate is not sharp at this scale.}
        The WKB outer-tail estimate uses a one-dimensional
        tunneling model.
        For $k = O(c^{1/3})$, the PSWF is
        not well approximated by a 1D Schr\"odinger
        eigenfunction; multi-dimensional effects
        or exact ODE analysis might give a
        larger lower bound.

  \item \textbf{Phase-space concentration vs.\ operator range.}
        $\pswf{k}$ for $k \leq N$ has Fourier support
        essentially in $[-c, c]$ and spatial support
        essentially in $[-1, 1]$.
        For $k = O(c^{1/3})$, the phase-space
        volume $4c \cdot 2 = 8c$ greatly exceeds $k$,
        meaning these modes \emph{underfill}
        their phase-space box.
        This underfilling could cause non-trivial
        leakage at polynomial scale.
\end{enumerate}

\subsection{Scenario (B): B-strong is false}

If $\varepsilon_N \ll c^{-1/2}$, then B-strong fails,
and the Paper~VI Galerkin estimate requires
a fundamentally different proof strategy.
In this case:
\begin{itemize}[nosep]
  \item The Galerkin matrix $G_N$ may not be
        invertible with $\|G_N^{-1}\| = O(c^{1/2})$.
  \item The error bound of Paper~VI, Theorem~1.1
        is not derivable via the B-strong route.
  \item Alternative: direct coercivity of the
        full Galerkin form without eigenvalue
        factorisation may still hold, but
        requires a different proof.
\end{itemize}

%% ============================================================
\section{Open Problems and Programme Decision Point}
\label{sec:decision}
%% ============================================================

\begin{openproblem}[Fredholm determinant precision]
\label{op:fredholm_precision}
Determine the asymptotic expansion of
$\log \Dc(z) = \log \det(I - z\Tc)$
for $z \in (0,1)$, $z = 1 - sc^{-1/2}$, to precision
$O(c^{1/3})$, explicitly separating the contribution
$\sum_{k=0}^N \log(1 - z\lam{k})$ from the bulk.
\end{openproblem}

\begin{openproblem}[$D_c(1)$ and Galerkin product]
\label{op:dc1}
Determine the asymptotic order of
$\Dc(1) = \prod_{k \geq 0} \varepsilon_k^{(c)}$
as $c \to \infty$.
If $\Dc(1)$ has polynomial decay (i.e.
$\log |\Dc(1)| = O(\log c)$),
then by AM--GM, at least one Galerkin eigenvalue
residual $\varepsilon_k$ is polynomial.
\emph{This would confirm scenario~(A).}
If $\log |\Dc(1)| = -\Omega(c^{2/3})$, scenario~(B)
becomes more plausible.
\end{openproblem}

\begin{openproblem}[Eigenvalue repulsion at Galerkin scale]
\label{op:repulsion}
Using the Hadamard--Weyl formula or the kernel
determinant structure of the prolate operator,
estimate the spacing $\lam{N} - \lam{N+1}$
for $N = O(c^{1/3})$.
If this spacing is $\Omega(c^{-1/2})$,
it supports polynomial order of $\varepsilon_N$
via a level-repulsion argument.
\end{openproblem}

\begin{remark}[Programme decision point]
\label{rem:decision}
The programme now faces a genuine bifurcation:
\begin{center}
\renewcommand{\arraystretch}{1.5}
\begin{tabular}{p{2.5cm} p{5cm} p{5cm}}
\toprule
& \textbf{Scenario (A): B-strong true}
& \textbf{Scenario (B): B-strong false} \\
\midrule
Mechanism
& Polynomial $\varepsilon_N$ via cancellation
& WKB bound is sharp; $\varepsilon_N \ll c^{-1/2}$ \\
Consequence
& Programme closes; Paper VI Thm 1.1 proven
& New proof strategy needed for coercivity \\
Next step
& Op.~Prob.~\ref{op:dc1}: control $D_c(1)$
& Direct coercivity of Galerkin form without $P_{kl}$ \\
Difficulty
& Fredholm precision; cancellation mechanism
& Full Paper VI restructuring \\
\bottomrule
\end{tabular}
\end{center}
\end{remark}

\medskip\noindent
\textbf{Recommended next step.}
Open Problem~\ref{op:dc1} is the minimal
numerically testable statement that distinguishes
the two scenarios:
compute $\log |\Dc(1)|$ numerically via
Bornemann-style Gauss--Legendre discretisation
(see \texttt{numerics/fredholm\_det\_scaling.py})
and check the growth rate as $c \to \infty$.

\begin{remark}[Numerical evidence for bulk dominance]
\label{rem:numerics}
A Bornemann-style discretisation of $\Tc$ on
$[-1,1]$ (Gauss--Legendre, $m = 6c$ nodes)
was computed for $c \in \{5, 8, 10, 15, 20, 30, 40, 50, 60, 75, 100\}$.
The results are recorded in
\texttt{numerics/fredholm\_det\_results.csv}.
The observed scaling is:
\[
  \log |\Dc(1)| \approx -18.4 \cdot c,
  \qquad c \in [5, 100],
\]
with effective local exponent
$\frac{d \log |\log \Dc|}{d \log c} \approx 1.18$
(stabilising from above as $c$ increases).

\medskip\noindent
\textbf{Identifiability limit.}
On the range $c \leq 100$, the two models
$\alpha c^p$ (best fit: $p \approx 1.24$, $R^2 = 0.9957$)
and $\gamma c \log c$ ($R^2 = 0.9975$)
are \emph{not distinguishable} from residual structure:
both exhibit a systematic U-shaped residual pattern
of comparable amplitude.
The effective exponent $p \approx 1 + 1/\log c_{\mathrm{eff}}$
is consistent with a $c \log c$ correction
evaluated at intermediate $c$.
Separation of the two models requires $c \gtrsim 500$,
which lies outside the current numerical range.

\medskip\noindent
\textbf{Conservative conclusion.}
The data confirm that $\log |\Dc(1)|$ grows at
least linearly in $c$, i.e.\
$\log |\Dc(1)| = -\Omega(c)$.
In particular, $\Dc(1)$ does \emph{not} have
polynomial order, ruling out the AM--GM argument
for a polynomial $\varepsilon_k$ from global
determinant data alone.
This is consistent with scenario~(B),
but does not exclude scenario~(A) via
a different mechanism.
The bulk dominance of $\log |\Dc(1)|$ confirms
that $\Dc(1)$ carries no accessible Galerkin-scale
signal on the current numerical range.
\end{remark}

%% ============================================================
\section{Microlocal Interpretation}
\label{sec:microlocal}
%% ============================================================

\subsection{Effective spectral density}

The preceding analysis --- both theoretical
(Proposition~\ref{prop:logdet}) and numerical
(Remark~\ref{rem:numerics}) --- establishes that
$\log |\Dc(1)|$ is governed by a continuous
spectral density, not by individual eigenvalues.
We make this precise.

\begin{definition}[Effective spectral action]
\label{def:phi}
Define the \emph{effective spectral action} of $\Tc$ by:
\[
  \Phi(c) := -\log |\Dc(1)|
  = \sum_{k=0}^\infty (-\log \eps{k})
  = \sum_{k=0}^\infty (-\log(1 - \lam{k})).
\]
By Proposition~\ref{prop:fredholm_basic}(1),
$\Phi(c) \geq 0$.
Numerically, $\Phi(c) \sim 18.4\, c$ for $c \in [5,100]$
(Remark~\ref{rem:numerics}).
\end{definition}

\begin{proposition}[Decomposition of $\Phi(c)$]
\label{prop:phi_decomp}
Write the sum defining $\Phi(c)$ as three contributions:
\begin{align*}
  \Phi(c)
  &= \underbrace{\sum_{k=0}^{N} (-\log \eps{k})}_{\Phi_{\mathrm{Gal}}(c)}
  + \underbrace{\sum_{k=N+1}^{N_{\mathrm{Sh}} - Cc^{1/3}}
    (-\log \eps{k})}_{\Phi_{\mathrm{bulk}}(c)}
  + \underbrace{\sum_{|k - N_{\mathrm{Sh}}| \leq Cc^{1/3}}
    (-\log \eps{k})}_{\Phi_{\mathrm{Airy}}(c)}.
\end{align*}
The orders of magnitude are:
\begin{center}
\renewcommand{\arraystretch}{1.4}
\begin{tabular}{llll}
\toprule
Term & \# modes & per-mode contrib. & total \\
\midrule
$\Phi_{\mathrm{Gal}}$ & $N \sim c^{1/3}$
  & $-\log\eps{k} \gtrsim \pi c$ (WKB)
  & $\Omega(c^{4/3})$ \\
$\Phi_{\mathrm{bulk}}$ & $\sim N_{\mathrm{Sh}} \sim c$
  & $-\log\eps{k} \sim \pi c \cdot (k/c)$ (LP)
  & $O(c^2)$ \\
$\Phi_{\mathrm{Airy}}$ & $O(c^{1/3})$
  & $O(\log c)$ (Airy CDF)
  & $O(c^{1/3} \log c)$ \\
\bottomrule
\end{tabular}
\end{center}
\end{proposition}

\begin{remark}[Airy is subleading]
\label{rem:airy_subleading}
Proposition~\ref{prop:phi_decomp} shows that
$\Phi_{\mathrm{Airy}}(c) = O(c^{1/3} \log c)$
is strictly subleading relative to
$\Phi_{\mathrm{bulk}}(c) = O(c^2)$
and even to $\Phi_{\mathrm{Gal}}(c) = \Omega(c^{4/3})$.

\medskip\noindent
This explains why the numerically observed effective
exponent $\approx 1.18 > 1$ is \emph{not} a signature
of Airy-edge dominance:
the Airy contribution is too small to be visible
at the precision accessible in the range $c \leq 100$.
The dominant signal is the bulk integral
$\Phi_{\mathrm{bulk}}$, with the Galerkin sector
$\Phi_{\mathrm{Gal}}$ contributing at an
intermediate scale.

\medskip\noindent
The non-linear coupling question of Paper~XI ---
whether the Airy transition generates a
spectral rigidity effect that forces $\varepsilon_N$
to be polynomial --- is therefore \emph{invisible}
in $\Phi(c)$ at any accessible numerical precision.
It is a subleading-of-subleading effect.
\end{remark}

\subsection{The transition proposition}

\begin{proposition}[Closure of the numerical phase]
\label{prop:closure}
The following has been established:
\begin{enumerate}[nosep]
  \item $\Phi(c) = -\Omega(c)$
        (numerically confirmed, $c \leq 100$).
  \item $\Phi_{\mathrm{Airy}}(c) = O(c^{1/3}\log c)$
        is subleading (analytical, from Proposition~\ref{prop:phi_decomp}).
  \item The Galerkin sector $\Phi_{\mathrm{Gal}}(c)$
        contributes at scale $\Omega(c^{4/3})$,
        but is masked by the bulk in $\Phi(c)$.
  \item No accessible numerical or analytical method
        extracts $\varepsilon_N$ from $\Phi(c)$
        at the required precision
        (Proposition~\ref{prop:precision}).
\end{enumerate}
\noindent
The present analysis therefore determines the
leading-order growth of $\log|\Dc(1)|$,
but does \emph{not} resolve whether the subleading
structure is governed by a logarithmic correction
or a higher-order spectral rigidity effect.

\medskip\noindent
The resolution of the dichotomy (A) vs.\ (B)
requires a method that bypasses $\Phi(c)$ entirely
and addresses $\varepsilon_N$ directly ---
either via a non-linear Fredholm symbol calculus,
or via the direct operator-theoretic approach
(coercivity of $G_N$ without $P_{kl}$ factorisation).
\end{proposition}

\begin{remark}[What comes next]
Proposition~\ref{prop:closure} marks the
structural transition of the programme:
\begin{itemize}[nosep]
  \item The \emph{numerical phase} is complete.
  \item The \emph{operator-theoretic phase} begins.
  \item The single remaining question is:
        \emph{can one prove or disprove
        $\varepsilon_N \geq Cc^{-1/2}$
        without going through $\Dc(1)$?}
\end{itemize}
This is the content of Paper~XII.
\end{remark}

%% ============================================================
\begin{thebibliography}{9}

\bibitem{paper9}
  \textit{Paper~IX: B-strong as a Spectral Gap},
  \texttt{paper9\_superconcentration.tex},
  prolate-gram-coercivity programme, 2026.

\bibitem{paper10}
  \textit{Paper~X: Spectral Gap Lower Bound},
  \texttt{paper10\_coercivity\_gap.tex},
  prolate-gram-coercivity programme, 2026.

\bibitem{LP1961II}
  H.~J.~Landau and H.~O.~Pollak,
  ``Prolate Spheroidal Wave Functions,
  Fourier Analysis and Uncertainty~II,''
  \textit{Bell Syst.\ Tech.\ J.}, 40 (1961), 65--84.

\bibitem{ORX2013}
  A.~Osipov, V.~Rokhlin, H.~Xiao,
  \textit{Prolate Spheroidal Wave Functions of Order Zero},
  Springer, 2013.

\bibitem{Simon2005}
  B.~Simon,
  \textit{Trace Ideals and Their Applications},
  2nd ed., AMS, 2005.

\bibitem{Bornemann2010}
  F.~Bornemann,
  ``On the numerical evaluation of Fredholm determinants,''
  \textit{Math.\ Comp.}, 79 (2010), 871--915.

\end{thebibliography}

\end{document}
